% ============================================================
%  SECCIÓN — PROCESO CONSTRUCTIVO Y DIAGRAMA DE GANTT (ARQUETA)
% ============================================================
\section{Proceso constructivo y programación de obras}
\label{sec:arqueta_proceso}

% -----------------------------------------------------------
\subsection{Descripción del proceso}
% -----------------------------------------------------------

La arqueta de fangos mezcla ($4{,}60 \times 3{,}60$~m en planta,
$H_\text{ext}=4{,}10$~m) se ejecuta como una
\textbf{estructura de hormigón armado in situ} en desmonte abierto,
sin entibación ni contención provisional, dado el reducido volumen de
excavación ($V_\text{exc}\approx 60$~m³) y la buena cohesión del
estrato~E-1a (arena limosa).
La cimentación es una losa superficial HA-30 a $D_f = 1{,}80$~m.

No se requiere pilotaje: el asiento total estimado
($s_\text{tot}=2{,}51$~cm $<$ 5~cm admisible) y la capacidad portante
(FS$=15{,}2$) cumplen holgadamente con criterio de cimentación directa.

Las etapas principales son:

\begin{enumerate}[label=\textbf{\arabic*.}]
  \item \textbf{Replanteo y trabajos previos.}
    Marcado del perímetro de excavación con holgura de 0,70~m por lado,
    señalización de obra y accesos para maquinaria.
  \item \textbf{Excavación y hormigón de limpieza.}
    Excavación a cielo abierto con retroexcavadora hasta la cota
    $-(D_f + 0{,}10)~\text{m} = -1{,}90$~m.
    Vertido de HM-15 ($e=0{,}10$~m) sobre el fondo saneado.
  \item \textbf{Losa de cimentación.}
    Ferrallado con B~500~SD (recubrimiento nominal 45~mm) y encofrado
    perimetral. Hormigonado con HA-30/B/20/IIb+Qc.
    Curado mínimo de \textbf{7~días} antes de alzar las paredes.
  \item \textbf{Paredes perimetrales.}
    Ferrallado de las cuatro paredes. Encofrado a dos caras (interior
    y exterior). Hormigonado con HA-30 y curado mínimo de
    \textbf{14~días} para garantizar resistencia suficiente antes
    de cargar la tapa.
  \item \textbf{Tapa superior.}
    Montaje de cimbra y encofrado inferior. Ferrallado y hormigonado
    de la tapa ($e=0{,}30$~m). Curado mínimo de 7~días.
  \item \textbf{Desencofrado e impermeabilización.}
    Retirada del encofrado interior y exterior. Aplicación del sistema
    de impermeabilización en extradós: lámina HDPE $\geq 2$~mm
    + geotextil 150~g/m² en todas las superficies enterradas
    ($\approx 41$~m²).
  \item \textbf{Relleno, compactación y prueba de estanqueidad.}
    Relleno del trasdós con material seleccionado en tongadas de
    30~cm, compactadas al 95\,\% Proctor modificado.
    Prueba hidráulica de estanqueidad conforme a normativa
    antes de la puesta en servicio.
\end{enumerate}

\begin{notabox}[Comparativa de plazos: arqueta vs.\ muro]
  La arqueta ($4{,}60 \times 3{,}60$~m en planta,
  $\approx 26$~m³ de HA-30 estructural)
  requiere un plazo total de \textbf{12~semanas} frente a las
  \textbf{31~semanas} del muro de contención
  ($L = 150$~ml, $\approx 1\,190$~m³ de HA-30).
  Los rendimientos adoptados son los mismos en ambas obras; la
  diferencia de plazo refleja la diferencia de volumen de trabajo:
  el muro tiene $\approx 45\times$ más hormigón estructural.
  Contribuyen además a reducir el plazo de la arqueta la ausencia
  de pilotaje (cimentación superficial) y el encofrado convencional
  sin sistema trepante.
\end{notabox}

% -----------------------------------------------------------
\subsection{Cálculo de duraciones}
\label{subsec:duraciones_arqueta}
% -----------------------------------------------------------

Las duraciones se obtienen dividiendo la medición de cada actividad
entre el rendimiento de la cuadrilla correspondiente y redondeando
al número entero de semanas por exceso.
Se aplican los mismos rendimientos de referencia que en el muro
(tabla~\ref{tab:actividades_arqueta}).

\begin{table}[H]
  \centering
  \caption{Cálculo de duraciones por actividad --- arqueta de fangos mezcla}
  \label{tab:actividades_arqueta}
  \resizebox{\textwidth}{!}{%
  \begin{tabular}{@{}clccccc@{}}
    \toprule
    \textbf{ID} & \textbf{Actividad}
      & \textbf{Medición} & \textbf{Rendimiento}
      & \textbf{Días} & \textbf{Semanas} & \textbf{Pred.} \\
    \midrule
    A & Replanteo y trabajos previos
      & ---
      & ---
      & 3 & 1 & --- \\
    B & Excavación + HM-15 de limpieza
      & 60~m³ / 1,7~m³
      & 300 / 30~m³/día
      & 1 & 1 & A \\
    C & Ferrallado y encofrado de losa
      & 400~kg / 5~m²
      & 1\,000~kg/día\ /\ 40~m²/día
      & 1 & 1 & B \\
    D & Hormigonado y curado de losa
      & 5,0~m³
      & 30~m³/h $+$ 7~d~cur.
      & 8 & 2 & C \\
    E & Ferrallado y encofrado de paredes
      & 1\,280~kg / 106~m²
      & 1\,000~kg/día\ /\ 40~m²/día
      & 4 & 2 & D* \\
    F & Hormigonado y curado de paredes
      & 16,0~m³
      & 30~m³/h $+$ 14~d~cur.
      & 15 & 2 & E \\
    G & Ferrallado y encofrado de tapa
      & 400~kg / 17~m²
      & 1\,000~kg/día\ /\ 40~m²/día
      & 1 & 1 & F \\
    H & Hormigonado y curado de tapa
      & 5,0~m³
      & 30~m³/h $+$ 7~d~cur.
      & 8 & 2 & G \\
    I & Desencofrado e impermeabilización
      & 106~m² / 41~m²
      & 40 / 200~m²/día
      & 3 & 1 & H \\
    J & Relleno, compactación y prueba
      & 47~m³
      & 150~m³/día
      & 2 & 1 & I\dag \\
    \midrule
    \multicolumn{5}{@{}l}{\textbf{PLAZO TOTAL (ruta crítica: A-B-C-D-E-F-G-H-I-J)}}
      & \textbf{12} & \\
    \bottomrule
  \end{tabular}}
\end{table}

\begin{notabox}[Notas sobre predecesoras y solapamientos]
  \begin{itemize}[leftmargin=1.2em, itemsep=2pt]
    \item[*] \textbf{E inicia con 1~semana de solape sobre D.}
      El ferrallado de paredes puede comenzar durante la
      última semana de curado de la losa (la resistencia
      para recibir las cargas de montaje del encofrado de
      paredes se alcanza antes de los 7~días completos).
    \item[\dag] \textbf{J se solapa totalmente con el final de I.}
      El relleno de trasdós y la prueba de estanqueidad
      se ejecutan en paralelo con el desencofrado de la cara
      exterior, dado el reducido volumen de trabajo de ambas
      actividades ($< 2$~días cada una).
  \end{itemize}
\end{notabox}

% -----------------------------------------------------------
\subsection{Diagrama de Gantt}
\label{subsec:gantt_arqueta}
% -----------------------------------------------------------

La figura~\ref{fig:gantt_arqueta} recoge la programación sobre un
horizonte de \textbf{12~semanas} ($\approx 3$~meses), con inicio en
la semana~1 y finalización en la semana~12.

\begin{figure}[H]
\centering
\resizebox{\textwidth}{!}{%
\begin{ganttchart}[
    hgrid,
    vgrid={*{1}{dotted}},
    x unit=1.10cm,
    y unit title=0.60cm,
    y unit chart=0.55cm,
    title label font=\bfseries\footnotesize,
    bar label font=\footnotesize,
    bar/.style={fill=BurntOrange!45, rounded corners=2pt},
    bar label node/.append style={text width=4.4cm}
  ]{1}{12}
  \gantttitle{Programación de obras --- arqueta de fangos mezcla (semanas)}{12} \\
  \gantttitlelist{1,...,12}{1} \\
  \ganttbar{A -- Replanteo}{1}{1} \\
  \ganttbar{B -- Excavaci\'{o}n + HM-15}{2}{2} \\
  \ganttbar{C -- Ferr.\ + enc.\ losa}{3}{3} \\
  \ganttbar{D -- Horm.\ + curado losa}{4}{5} \\
  \ganttbar{E -- Ferr.\ + enc.\ paredes}{5}{6} \\
  \ganttbar{F -- Horm.\ + curado paredes}{7}{8} \\
  \ganttbar{G -- Ferr.\ + enc.\ tapa}{9}{9} \\
  \ganttbar{H -- Horm.\ + curado tapa}{10}{11} \\
  \ganttbar{I -- Desencof.\ + impermeab.}{12}{12} \\
  \ganttbar{J -- Relleno + acabados}{12}{12}
\end{ganttchart}%
}
\caption{Diagrama de Gantt del proceso constructivo de la arqueta de fangos mezcla
         (plazo total: 12~semanas $\approx$ 3~meses).
         Actividades I y J se ejecutan en paralelo en la semana~12.}
\label{fig:gantt_arqueta}
\end{figure}
